% ------------------------------------------------------------
% SEDONA SPINE: COMPLETE PREAMBLE
% Document: Defensive Publication – Prime-Indexed Operator Calculus
% ------------------------------------------------------------

% ---------- Document Class ----------
\documentclass[12pt]{article}

% ---------- Page Geometry ----------
\usepackage[letterpaper, margin=1in, headheight=15pt, includehead]{geometry}

% ---------- Font Encoding ----------
\usepackage[utf8]{inputenc}
\usepackage[T1]{fontenc}
\usepackage{textcomp}
\usepackage{lmodern}

% ---------- Language and Hyphenation ----------
\usepackage[english]{babel}

% ---------- Graphics and Colors ----------
\usepackage{graphicx}
\usepackage{xcolor}
\usepackage{epstopdf}
\usepackage{tikz}
\usetikzlibrary{positioning, arrows, shapes, calc, fit, backgrounds}

% ---------- Mathematics ----------
\usepackage{amsmath, amssymb, amsthm, mathtools}
\usepackage{bm}
\usepackage{dsfont}
\usepackage{bbm}
\usepackage{stmaryrd}       % for \llbracket, \rrbracket

% ---------- Tables ----------
\usepackage{array}
\usepackage{booktabs}
\usepackage{longtable}
\usepackage{multirow}
\usepackage{colortbl}

% ---------- Listings and Code ----------
\usepackage{listings}
\usepackage{verbatim}
\usepackage{moreverb}

% ---------- Hyperlinks and Cross-referencing ----------
\usepackage{hyperref}
\hypersetup{
    colorlinks=true,
    linkcolor=blue,
    urlcolor=blue,
    citecolor=blue,
    pdftitle={Sedona Spine: Prime-Indexed Operator Calculus},
    pdfauthor={Ryan and Contributors},
    pdfsubject={Defensive Publication}
}
\usepackage{cleveref}

% ---------- Captions and Floats ----------
\usepackage{caption}
\usepackage{subcaption}
\usepackage{float}

% ---------- Headers and Footers ----------
\usepackage{fancyhdr}
\pagestyle{fancy}
\fancyhf{}
\fancyhead[L]{\textit{Sedona Spine Defensive Publication}}
\fancyhead[R]{\thepage}
\renewcommand{\headrulewidth}{0.4pt}

% ---------- Custom Commands and Environments ----------

% ---- Mathematical Operators ----
\newcommand{\op}[1]{\operatorname{#1}}
\DeclareMathOperator{\Tr}{Tr}
\DeclareMathOperator{\diag}{diag}
\DeclareMathOperator{\Spec}{Spec}
\DeclareMathOperator{\Ker}{Ker}
\DeclareMathOperator{\Span}{Span}
\DeclareMathOperator{\Proj}{Proj}
\DeclareMathOperator{\Hash}{Hash}
\DeclareMathOperator{\SHA}{SHA}

% ---- Norms and Brackets ----
\providecommand{\norm}[1]{\left\lVert#1\right\rVert}
\providecommand{\abs}[1]{\left\lvert#1\right\rvert}
\providecommand{\inner}[1]{\left\langle#1\right\rangle}
\providecommand{\bra}[1]{\left\langle#1\right\rvert}
\providecommand{\ket}[1]{\left\lvert#1\right\rangle}
\providecommand{\braket}[2]{\left\langle#1 \middle| #2\right\rangle}

% ---- Sets and Spaces ----
\newcommand{\R}{\mathbb{R}}
\newcommand{\C}{\mathbb{C}}
\newcommand{\Z}{\mathbb{Z}}
\newcommand{\N}{\mathbb{N}}
\newcommand{\Q}{\mathbb{Q}}
\newcommand{\F}{\mathbb{F}}
\newcommand{\M}{\mathcal{M}}
\newcommand{\Hilb}{\mathcal{H}}
\newcommand{\Fock}{\mathcal{F}}

% ---- Operators from the Framework ----
\providecommand{\LambdaM}{\Lambda_m^{\text{op}}}
\providecommand{\XiOp}{\Xi}
\providecommand{\Zeno}{Z}
\providecommand{\PiM}{\Pi_M}
\providecommand{\nablaD}{\nabla D}
\providecommand{\rhoWarn}{\rho_{\text{warn}}}
\providecommand{\rhoStar}{\rho^*}
\providecommand{\rhoCrit}{\rho_{\text{crit}}}

% ---- Custom Math Operators for Zeta and Multiplicity ----
\DeclareMathOperator{\D}{\mathcal{D}}
\DeclareMathOperator{\ProjZeta}{Proj_\zeta}
\DeclareMathOperator{\PWEH}{PWEH}

% ---- Theorems, Definitions, Lemmas, Corollaries ----
\theoremstyle{plain}
\newtheorem{theorem}{Theorem}[section]
\newtheorem{lemma}[theorem]{Lemma}
\newtheorem{corollary}[theorem]{Corollary}
\newtheorem{proposition}[theorem]{Proposition}
\newtheorem{conjecture}[theorem]{Conjecture}

\theoremstyle{definition}
\newtheorem{definition}[theorem]{Definition}
\newtheorem{example}[theorem]{Example}
\newtheorem{remark}[theorem]{Remark}
\newtheorem{note}[theorem]{Note}

\theoremstyle{remark}
\newtheorem{acknowledgement}[theorem]{Acknowledgement}

% ---- Proof environment with QED ----
\renewenvironment{proof}{\paragraph{\textit{Proof}.}}{\hfill$\square$}

% ---- Code Listings Style ----
\lstset{
    basicstyle=\ttfamily\small,
    keywordstyle=\color{blue}\bfseries,
    commentstyle=\color{green!60!black},
    stringstyle=\color{red!70!black},
    numbers=left,
    numberstyle=\tiny\color{gray},
    breaklines=true,
    frame=single,
    captionpos=b,
    tabsize=4,
    showspaces=false,
    showstringspaces=false,
    showtabs=false,
    xleftmargin=2em,
    escapeinside={(*}{*)}
}

% ---- Custom Listing Languages ----
\lstdefinelanguage{Rust}{
    keywords={as, break, const, continue, crate, else, enum, extern, false, fn, for, if, impl, in, let, loop, match, mod, move, mut, pub, ref, return, self, Self, static, struct, super, trait, true, type, unsafe, use, where, while, async, await, dyn},
    keywordstyle=\color{blue}\bfseries,
    comment=[l]{//},
    morecomment=[s]{/*}{*/},
    string=[b]",
    string=[b]',
    sensitive=true
}

\lstdefinelanguage{Lean}{
    keywords={def, theorem, lemma, corollary, proof, variable, structure, import, by, intro, unfold, have, exact, apply, auto, rewrite, simp, induction, cases, split, use, exists, assume, if, then, else, forall, in, let, show},
    keywordstyle=\color{blue}\bfseries,
    comment=[l]{--},
    comment=[s]{/-}{-/},
    string=[b]",
    sensitive=true
}

% ---- Additional Useful Macros ----
\newcommand{\ceil}[1]{\left\lceil#1\right\rceil}
\newcommand{\floor}[1]{\left\lfloor#1\right\rfloor}
\newcommand{\eps}{\varepsilon}
\newcommand{\vect}[1]{\mathbf{#1}}
\newcommand{\mat}[1]{\mathbf{#1}}
\newcommand{\partialderiv}[2]{\frac{\partial #1}{\partial #2}}

% ---- Appendix and Section Numbering ----
\usepackage{titlesec}
\titleformat{\section}{\bfseries\Large}{\thesection}{1em}{}
\titleformat{\subsection}{\bfseries\large}{\thesubsection}{1em}{}
\titleformat{\subsubsection}{\bfseries\normalsize}{\thesubsubsection}{1em}{}

% ---- Enable PDF Metadata (optional) ----

\title{\bfseries Automorphic Transformers with Certified Graph Energetics and\\ an Application to Predictive Maintenance: A Defensive Publication \\ \Large\textit{In Legacy of Mubeen Mohammad}}
\author{Citizen Gardens \\ The Prime Materia Commons}
\date{\today}

\begin{document}

\maketitle

\begin{abstract}
We disclose a family of deeply structured neural architectures and training protocols in which every inductive bias -- symmetry, spectral regularity, Lipschitz stability, and data quality -- is encoded as a testable certificate, a constrained projection, or a statistically grounded hypothesis test. The core is an \textbf{automorphic transformer} whose attention pattern is governed by finite affine groups over prime fields, with a mask derived from the quadratic character interpreted as a $2$-cocycle twist. The system enforces invariance of class functions under these group actions (CSL loss), unitarizes attention matrices via Sinkhorn–Cayley/exp paths, tests eigenphase distributions against Sato--Tate laws, and imposes an $\ell_1$-budgeted projection with dual-gap certificates. A subsequent \textbf{AU--Triad} extension constructs a symmetric certification graph from the bistochastic attention, adds a Laplacian-based heat equalization step with a contraction certificate, and composites all diagnostics into a single per‑step audit tuple governed by a tiered gate policy. We then transpose the entire certified learning engine into the domain of industrial predictive maintenance as an \textbf{Integrated Machine Health Pipeline (IMHP)}: an edge‑to‑cloud agent that ingests MQTT telemetry from ESP32 sensors, validates data quality, extracts condition features, and runs anomaly detection and remaining useful life (RUL) forecasting inside a certifiable machine digital twin. The modular orchestration (ingest--validate--condition--infer--decide--alert) folds data‑trustworthiness scoring into the certification framework, so that every maintenance recommendation carries both model‑level and data‑provenance guarantees. This disclosure establishes prior art for the entire stack: the finite automorphic bias, the cocycle mask, the certified spectral/graph loops, and their concrete embedding in a predictive‑maintenance agent.
\end{abstract}

\section{Introduction}
Modern deep learning systems often rely on heuristics and post‑hoc explanations. We describe a complete, mathematically auditable architecture where every claimed property -- symmetry, spectral regularity, Lipschitz stability, and data integrity -- is associated with an explicit certificate, loss, or hypothesis test. The design begins with a \textbf{transformer whose attention structure is tied to a concrete finite group}, continues through \textbf{certified spectral and graph‑energetic audits}, and finally is deployed as a \textbf{machine health digital twin} that monitors industrial assets via MQTT/ESP32/Node‑RED.

This document serves as a defensive publication to establish prior art for:
\begin{enumerate}
    \item Attention mechanisms that use additive masks derived from quadratic characters of finite fields (Legendre symbol) interpreted as $2$-cocycles of the action groupoid of $AGL(1,p)$ on token indices.
    \item Training objectives that enforce invariance of class functions under finite group actions (CSL loss).
    \item Unitarization of attention matrices via Sinkhorn iterations followed by exponential or Cayley maps, with eigenphase statistics tested against Sato--Tate distributions.
    \item Graph‑theoretic certification of attention: symmetrisation to a Laplacian, Dirichlet energy monitoring, and a certified heat‑flow equalization step with contraction guarantee.
    \item Per‑step composite audit tuples and tiered gate policies that decide acceptance, rollback, or abort.
    \item The integration of the above into an edge‑to‑cloud predictive‑maintenance agent (IMHP) with a dedicated data‑quality gating layer.
\end{enumerate}

\section{Mathematical Preliminaries}
\label{sec:prelim}

\subsection{Finite Fields and CRT Embedding}
Let $p_1, p_2$ be distinct odd primes such that $p_1 p_2 \ge N$, where $N$ is the number of tokens (e.g.\ sensor windows). Each token index $i \in [0, N-1]$ is mapped via the Chinese Remainder Theorem (CRT) to a pair $(i \bmod p_1,\; i \bmod p_2) \in \mathbb{F}_{p_1} \times \mathbb{F}_{p_2}$.

\subsection{The Group $AGL(1,p)$}
For a prime $p$, the affine group $AGL(1,p)$ consists of transformations $x \mapsto a x + b$ with $a \in \mathbb{F}_p^\times$, $b \in \mathbb{F}_p$. Its order is $p(p-1)$. The group acts on the set of indices in $\mathbb{F}_p$ and, via the CRT embedding, the direct product $G = AGL(1,p_1) \times AGL(1,p_2)$ acts on the token set. We consider the action of $G$ on attention matrices by simultaneous permutation of rows and columns: $A \mapsto P_g A P_g^\top$, where $P_g$ is the permutation matrix representing $g$.

\subsection{Quadratic Character and Cocycle Twist}
Let $\chi_p(x) = \left(\frac{x}{p}\right)$ be the Legendre symbol, with $\chi_p(0)=0$ by convention. For $x \neq 0$, $\chi_p(x) = +1$ if $x$ is a quadratic residue modulo $p$, and $-1$ otherwise. Consider the action groupoid $\mathcal{G}_p$ whose objects are indices in $\mathbb{F}_p$ and morphisms $i \to g(i)$ for $g \in AGL(1,p)$. The function
\[
c(i,j) = \chi_p(\varphi(i) - \varphi(j)),
\]
where $\varphi$ is the natural map from token index to $\mathbb{F}_p$, can be seen as a $\{\pm1\}$-valued $2$-cocycle on composable pairs of morphisms (restricting to edges where the difference is a quadratic residue defines a sub-groupoid; the cocycle weights the others). The \textbf{cocycle twist} perspective interprets the attention mask as a multiplicative twist of the groupoid convolution algebra.

\subsection{Class Functions and CSL Loss}
A class function $F_{\text{inv}}$ on the space of attention matrices is one that satisfies $F_{\text{inv}}(P_g A P_g^\top) = F_{\text{inv}}(A)$ for all $g \in G$. Examples include the sorted spectrum of a unitary derived from $A$, or the sorted row‑moments. The \textbf{CSL (Conjugacy‑class Specified Learning) loss} enforces this invariance:
\[
\mathcal{L}_{\text{CSL}} = \frac{1}{|S|} \sum_{g \in S} \Delta\big( F_{\text{inv}}(A), F_{\text{inv}}(P_g A P_g^\top) \big),
\]
where $S$ is a minibatch of group elements and $\Delta$ is a discrepancy metric (e.g.\ mean squared error of sorted eigenphases).

\section{Automorphic Attention with Legendre Cocycle Twists}
\label{sec:att}

The core attention operator for a single head is defined as follows. Let $Q,K \in \mathbb{R}^{N \times d}$ be query and key matrices. The raw logits are $L_{ij} = (Q K^\top)_{ij} / \sqrt{d}$. The \textbf{additive Legendre mask} is constructed from the quadratic character of the difference of CRT‑embedded indices:
\[
M_p[i,j] = \begin{cases}
+1 & \text{if } \chi_p(\varphi(i)-\varphi(j)) = +1,\\
-1 & \text{otherwise}.
\end{cases}
\]
The masked logits become
\[
\tilde{L}_{ij} = L_{ij} + \alpha M_p[i,j] - \beta (1 - M_p[i,j]),
\]
with $\alpha \ge 0$ and $\beta \gg 0$ penalising non‑residue edges. After softmax, one obtains a row‑stochastic matrix $\tilde{A}$. For multi‑prime settings the masks are applied independently and aggregated (e.g.\ elementwise AND).

\paragraph{Cocycle interpretation.} The term $\alpha M_p - \beta(1-M_p)$ is exactly the logarithm of a multiplicative factor $\exp( \alpha \chi_p(\Delta \varphi) )$ (with a large suppression for $\chi_p=-1$). This is the log of the $2$-cocycle twist on the groupoid convolution, turning the attention into a twisted convolution operator. The CSL loss then demands that the representation of this twisted algebra be invariant under conjugation by $G$.

\section{Spectral Unitarization and Sato--Tate Diagnostics}
\label{sec:unitary}

Row‑stochastic $\tilde{A}$ is stabilised by adding $\varepsilon$ proportional to the mean row sum, then mapped to a bistochastic matrix $B$ via a few ($\le 5$) Sinkhorn iterations:
\[
B^{(t+1)} = \operatorname{diag}(r^{(t)})^{-1} B^{(t)} \operatorname{diag}(c^{(t)})^{-1},
\]
where $r^{(t)}, c^{(t)}$ enforce row and column sums $=1$. 

From $B$ we obtain a unitary matrix $U$ either by the matrix exponential around the flat matrix $J = \frac{1}{N}\mathbf{1}\mathbf{1}^\top$:
\[
U = \exp\big(i \theta (B - J)\big),
\]
or via the Cayley transform $U = (I + iH)(I - iH)^{-1}$ with $H$ a suitable Hermitian matrix derived from $B$. These maps are permutation‑equivariant, so the multiset of eigenphases of $U$ is invariant under $G$.

We test the empirical distribution of eigenphases against the \textbf{Sato--Tate} measure $\frac{2}{\pi}\sin^2\theta \, d\theta$ (characteristic of random unitary matrices from certain families) using the Wasserstein‑2 distance (W2), Kolmogorov–Smirnov (KS), or Cramér–von Mises (CvM) statistics, with bootstrap confidence intervals (BCa). A run is considered \emph{ST‑compliant} if the median W2 and the CI width lie inside preregistered bands.

\section{Graph Energetics and Certified Heat Equalization (AU--Triad)}
\label{sec:au}

The AU--Triad extension (Measure–Compare–Equalize) constructs a \textbf{certification graph} from the bistochastic $B$:
\[
W = \frac{1}{2}(B + B^\top), \qquad D = \operatorname{diag}(W\mathbf{1}), \qquad L = D - W.
\]
$L$ is the symmetric graph Laplacian.

\subsection{Dirichlet Energy and Spectral Gap}
For a node signal $u \in \mathbb{R}^N$, the Dirichlet energy is $E[u] = \frac{1}{2} u^\top L u$. Let $0 = \lambda_1 \le \lambda_2 \le \dots \le \lambda_N$ be the eigenvalues of $L$. The \textbf{spectral gap} $\lambda_2$ controls mixing and smoothness.

\subsection{Certified Heat Equalization}
A heat‑flow step updates the signal:
\[
u^+ = u - \alpha L u, \quad 0 < \alpha < 2/\lambda_{\max}.
\]
The contraction ratio is bounded by $r_t \le |1 - \alpha \lambda_2|$ (for the component orthogonal to the constant vector). The step is \textbf{certified} if $r_t \le 1 - \alpha \lambda_2 \le 1 - \delta$ for a preregistered $\delta$. In the machine twin, $u$ is typically the channel‑mean of the token embeddings after the first automorphic block, but can be configured to other meaningful signals (e.g.\ fault logit projection).

\section{Lipschitz Envelopes and ACE Projection}
\label{sec:lipschitz}

\subsection{Softmax Jacobian Bound}
The Jacobian of the softmax operator $\sigma: \mathbb{R}^N \to \mathbb{R}^N$ satisfies $\|J_\sigma(x)\|_{1\to1} \le \frac12$. This gives a per‑layer Lipschitz bound (SoftmaxUB) that can be logged.

\subsection{Slope Upper Bound (SlopeUB)}
For a transformer stack, we compose per‑layer operator norms (attention linear maps, feed‑forward) with the SoftmaxUB to obtain a global $\ell_1$‑Lipschitz constant estimate, called \textbf{SlopeUB}. Training configurations must keep SlopeUB below a preregistered ceiling.

\subsection{ACE Projection}
After each gradient step, all trainable parameters are projected onto a weighted $\ell_1$‑ball:
\[
\Omega = \left\{ x : \sum_k \omega_k |x_k| \le \Lambda_m \right\}.
\]
The projection is the unique solution of a strictly convex KKT problem, computable via a separable shrinkage algorithm. The dual gap $\tau_{\text{ACE}}$ and the gap lower bound $\text{GapLB}$ are logged as optimality certificates.

\section{Composite Certification and Gate Policy}
\label{sec:gates}

Every training step produces a composite audit tuple:
\[
\mathcal{T} = \langle E[u], r_t, 1-\alpha\lambda_2, \lambda_2, \lambda_{\max}, \|U^*U-I\|_{1\to1}, \text{W2}, \text{KS}, \text{SlopeUB}, \tau_{\text{ACE}}, \text{gap}_{\text{dual}} \rangle.
\]
A tiered gate policy governs acceptance:

\paragraph{Burn‑in (first $K$ steps):}
\begin{itemize}[nosep]
    \item \textbf{Hard gates:} unitarity residual, ACE feasibility, ACE dual gap, Cayley safety, gross SlopeUB ceiling.
    \item \textbf{Soft gates:} Sato–Tate proxy W2/KS, heat contraction ratio, permutation invariance KS test. Violations are logged and penalised but do not cause rollback.
\end{itemize}

\paragraph{Post‑burn‑in:}
All soft gates become hard. A violation triggers rollback and retry (e.g.\ halve $\alpha$ on heat failure, fallback to exp unitarisation on Cayley failure). This yields a \textbf{certified step policy}.

\section{Application to Predictive Maintenance: IMHP Agent}
\label{sec:imhp}

The entire certified learning stack is deployed inside an \textbf{Integrated Machine Health Pipeline (IMHP)} -- a modular, edge‑to‑cloud agent for industrial predictive maintenance.

\subsection{System Architecture}
The IMHP is a supervisory agent with six specialised sub‑agents:
\begin{enumerate}[nosep]
    \item \textbf{Ingestion agent:} receives MQTT telemetry from ESP32 sensor nodes (vibration, temperature, current, RPM, etc.).
    \item \textbf{Validate agent:} enforces JSON schema, timestamp freshness, and unit consistency.
    \item \textbf{Condition agent (Quality + Features):} scores data trustworthiness (missingness, drift, outliers), then windows the cleaned streams and computes condition indicators (RMS, kurtosis, FFT band energies, thermal gradients).
    \item \textbf{Infer agent:} runs the automorphic machine twin -- anomaly detection, fault classification, RUL estimation, risk‑horizon labelling -- using the certified attention and AU–Triad loop.
    \item \textbf{Decide agent:} converts model outputs into maintenance recommendations (inspect, derate, lubricate, schedule) while incorporating prediction‑trust derived from data quality and gate compliance.
    \item \textbf{Alert agent:} publishes MQTT alerts and updates a Node‑RED dashboard.
\end{enumerate}

\subsection{Data‑Quality Gating}
A tiered trust policy mirrors the model gate policy:
\begin{itemize}[nosep]
    \item Short gaps in slow channels: forward fill with confidence penalty.
    \item Gaps in vibration windows: mark window degraded; suppress RUL updates.
    \item Long gaps/stale timestamps: downgrade asset status to ``telemetry unreliable''.
    \item Out‑of‑range values: quarantine as sensor fault unless corroborated.
    \item Repeated flatline: trigger sensor‑health alert.
\end{itemize}
Every feature vector carries a \texttt{quality\_score}; every prediction a \texttt{prediction\_trust}. These are appended to the audit tuple.

\subsection{MQTT Topic Structure}
\begin{verbatim}
imhp/{site}/{line}/{asset_id}/telemetry/{channel}
imhp/{site}/{line}/{asset_id}/health/anomaly
imhp/{site}/{line}/{asset_id}/health/rul
imhp/{site}/{line}/{asset_id}/health/certification
imhp/{site}/{line}/{asset_id}/alerts
\end{verbatim}

\subsection{Integration with the Certified Twin}
The automorphic twin runs inside the infer agent. Its per‑step audit tuple (Sec.~\ref{sec:gates}) is published on the \texttt{certification} topic, allowing operators to monitor not just predictions but the trustworthiness of the underlying model invariants.

\section{Code Snippets}
\label{sec:code}

The following Python‑style snippets illustrate key algorithms.

\subsection{Legendre Mask Construction}
\begin{lstlisting}[language=Python, caption={Building the additive Legendre mask as log of cocycle twist.}]
def legendre_mask(indices_1, indices_2, p, alpha=0.0, beta=20.0):
    # indices_1, indices_2: tensors of CRT values mod p
    diff = (indices_1.unsqueeze(1) - indices_2.unsqueeze(0)) % p
    # Legendre symbol: +1 for residues, -1 for non-residues, 0 for zero
    chi = torch.zeros_like(diff, dtype=torch.float)
    for x in range(1, p):
        chi[diff == x] = legendre_symbol(x, p)
    # mask: +1 for residue edges, 0 otherwise
    mask = (chi == 1).float()
    # additive logits: encourage +alpha on residues, penalise -beta on non-residues
    logit_bias = alpha * mask - beta * (1 - mask)
    return logit_bias
\end{lstlisting}

\subsection{Sinkhorn Bistochasticization}
\begin{lstlisting}[language=Python, caption={$\varepsilon$-Sinkhorn to obtain bistochastic $B$ from row-stochastic $\tilde{A}$.}]
def sinkhorn(A_tilde, eps=1e-6, iters=5):
    N = A_tilde.shape[0]
    B = A_tilde + eps / N
    for _ in range(iters):
        row_sum = B.sum(dim=1, keepdim=True)
        B = B / row_sum
        col_sum = B.sum(dim=0, keepdim=True)
        B = B / col_sum
    return B
\end{lstlisting}

\subsection{Heat Equalization with Contraction Certificate}
\begin{lstlisting}[language=Python, caption={Certified heat step on signal $u$ using Laplacian $L$.}]
def heat_step(u, L, alpha, lambda_max, lambda2):
    # 0 < alpha < 2/lambda_max
    u_new = u - alpha * (L @ u)
    contraction_ratio = torch.norm(u_new - u_new.mean()) / (torch.norm(u - u.mean()) + 1e-8)
    # certificate: ratio <= 1 - alpha * lambda2
    cert_passed = contraction_ratio <= (1 - alpha * lambda2 + 1e-5)
    return u_new, contraction_ratio, cert_passed
\end{lstlisting}

\subsection{Composite Gate Logic (sketch)}
\begin{lstlisting}[language=Python, caption={Tiered gate check with burn‑in.}]
def check_gates(audit, step, burn_in_steps, hard_list, soft_list):
    results = {}
    for gate_name, value, threshold in audit:
        if gate_name in hard_list or (step > burn_in_steps and gate_name in soft_list):
            results[gate_name] = (value <= threshold)
        else:
            results[gate_name] = True  # soft during burn-in
    return all(results.values()), results
\end{lstlisting}

\section{Discussion and Conclusion}
We have disclosed a complete, certificated machine learning pipeline whose mathematical foundations range from finite group actions and twisted groupoid algebras to graph Laplacian contraction and constrained convex optimisation. The architecture translates every design choice into a testable invariant: the Legendre mask encodes a $2$-cocycle twist, the CSL loss enforces class‑function invariance, the unitarization and Sato–Tate tests certify spectral regularity, the Laplacian heat step adds energetic contraction guarantees, and the ACE projection maintains parameter‑space Lipschitz budgets. All these certificates are aggregated into a single audit tuple and governed by a tiered gate policy.

The subsequent transposition to an edge‑to‑cloud predictive‑maintenance agent (IMHP) demonstrates practical deployability. By adding a data‑quality gating layer and folding it into the certificate tuple, the system ensures that maintenance decisions are not only model‑certified but also data‑provenance aware. This defensive publication establishes prior art for the entire stack, including the specific combination of automorphic attention, certified graph energetics, and the modular IMHP agent.

\vspace{0.5cm}
\noindent\textbf{Acknowledgment:} The authors thank the open‑source community for tools and discussions that made this synthesis possible.

\nocite{*}
\bibliographystyle{plain}
\bibliography{references}
\end{document}